%% ============================================================
%%  Multi-Customer, Multi-Item Inventory Management Under SLAs
%%  A Deep Controlled Learning Approach with Dynamic Allocation
%%  Enhanced Version with Detailed Experiments and DCL Framework
%% ============================================================

\documentclass[11pt, a4paper]{article}

%% ---- Packages ----
\usepackage[utf8]{inputenc}
\usepackage[T1]{fontenc}
\usepackage{lmodern}
\usepackage{pifont}
\usepackage[english]{babel}
\usepackage{amsmath, amssymb, amsthm}
\usepackage{mathtools}
\usepackage{booktabs}
\usepackage{array}
\usepackage{tabularx}
\usepackage{multirow}
\usepackage{graphicx}
\usepackage{xcolor}
\usepackage{hyperref}
\usepackage{natbib}
\usepackage{geometry}
\usepackage{setspace}
\usepackage{microtype}
\usepackage{enumitem}
\usepackage{caption}
\usepackage{subcaption}
\usepackage{float}
\usepackage{algorithm}
\usepackage{algpseudocode}
\usepackage{tikz}
\usepackage{longtable}
\usetikzlibrary{shapes.geometric, arrows.meta, positioning, calc}

%% ---- Page Layout ----
\geometry{
    a4paper,
    top=2.5cm, bottom=2.5cm,
    left=2.5cm, right=2.5cm
}
\onehalfspacing

%% ---- Hyperref Setup ----
\hypersetup{
    colorlinks=true,
    linkcolor=blue!60!black,
    citecolor=blue!60!black,
    urlcolor=blue!60!black,
    pdftitle={Multi-Customer Inventory Management Under Heterogeneous SLAs},
    pdfauthor={Eliz Payasli}
}

%% ---- Theorem Environments ----
\newtheorem{theorem}{Theorem}[section]
\newtheorem{lemma}[theorem]{Lemma}
\newtheorem{proposition}[theorem]{Proposition}
\newtheorem{corollary}[theorem]{Corollary}
\theoremstyle{definition}
\newtheorem{definition}[theorem]{Definition}
\newtheorem{example}[theorem]{Example}
\theoremstyle{remark}
\newtheorem{remark}[theorem]{Remark}

%% ---- Custom Commands ----
\newcommand{\IL}{\text{IL}}
\newcommand{\IP}{\text{IP}}
\newcommand{\AFR}{\text{AFR}}
\newcommand{\BO}{\text{BO}}
\newcommand{\Q}{\mathbf{Q}}
\newcommand{\S}{\mathbf{S}}
\newcommand{\D}{\mathbf{D}}
\newcommand{\E}{\mathbb{E}}
\newcommand{\Real}{\mathbb{R}}
\newcommand{\Nat}{\mathbb{N}}

%% ============================================================
\begin{document}

%% ---- Title Block ----
\begin{titlepage}
    \centering
    \vspace*{1.5cm}

    {\Large \textbf{Eindhoven University of Technology}\\[0.3em]
    Department of Industrial Engineering and Innovation Sciences\\
    Operations, Planning, Accounting \& Control}

    \vspace{3cm}

    {\LARGE \textbf{Dynamic Inventory Allocation\\[0.5em]
    in Multi-Customer Systems\\[0.5em]
    with Heterogeneous Service Level Agreements\\[0.5em]
    A Deep Controlled Learning Approach}}

    \vspace{2.5cm}

    {\large \textit{Research Paper --- Model Implementation}}

    \vspace{2cm}

    {\large
    \textbf{Author:} Eliz Payasli \\[0.7em]
    \textbf{Supervisors:} Willem van Jaarsveld, Tarkan Temizoz \\[0.3em]
    \textbf{Date:} September 2026 \\[0.3em]

    }

\end{titlepage}

%% ---- From here on, body text is left-aligned (ragged right) ----
\raggedright

\newpage
\section{Problem Description}

We consider a supplier that manages inventory at a single location and serves multiple customers under heterogeneous Service Level Agreements (SLAs). The supplier stocks a set of items $I$ (indexed $i \in I$, $|I| \geq 1$) and serves a set of customers $C$ (indexed $c \in C$, $|C| \geq 1$). Each customer-item pair has its own, generally heterogeneous, stochastic demand process, so different customers place very different demands on the same shared stock.

\begin{figure}[H]
    \centering
    \begin{tikzpicture}[
        node distance=1.8cm and 2.8cm,
        source/.style={rectangle, draw, thick, dashed, rounded corners, minimum width=3.4cm, minimum height=1.1cm, align=center, fill=gray!3},
        supplier/.style={rectangle, draw, thick, rounded corners, minimum width=3.4cm, minimum height=1.4cm, align=center, fill=gray!10},
        customer/.style={rectangle, draw, thick, rounded corners, minimum width=3.0cm, minimum height=1.2cm, align=center, fill=blue!5},
        arrow/.style={-{Latex[length=2.5mm]}, thick}
    ]
        \node[source] (Src) {External source};
        \node[supplier, below=of Src] (S) {Supplier\\[2pt] \footnotesize items $i \in I$, base-stock $S_i$};
        \node[customer, below left=of S] (C1) {Customer 1\\[2pt] \footnotesize $\beta_1$, demand rate $\lambda_{1,i}$};
        \node[customer, below=of S] (C2) {Customer 2\\[2pt] \footnotesize $\beta_2$, demand rate $\lambda_{2,i}$};
        \node[customer, below right=of S] (C3) {Customer $|C|$\\[2pt] \footnotesize $\beta_{|C|}$, demand rate $\lambda_{|C|,i}$};

        \draw[arrow] (Src) -- (S) node[midway, right, font=\scriptsize] {$Q_i(t) = \max(0, S_i - \text{IP}_i(t))$};
        \draw[arrow] (S) -- (C1) node[midway, above left, sloped, font=\scriptsize] {$A_{1,i}(t)$};
        \draw[arrow] (S) -- (C2) node[midway, right, font=\scriptsize] {$A_{2,i}(t)$};
        \draw[arrow] (S) -- (C3) node[midway, above right, sloped, font=\scriptsize] {$A_{|C|,i}(t)$};
    \end{tikzpicture}
    \caption{Two decisions govern the supplier's operation each period: (RQ1) an \emph{ordering} decision $Q_i(t)$ that replenishes stock from an external source under a static base-stock policy, and (RQ2) an \emph{allocation} decision $A_{c,i}(t)$ that dynamically splits available stock across customers $c \in C$ once demand is realized. Customers differ in demand rate $\lambda_{c,i}$ and in their SLA backorder allowance $\beta_c$ — a stricter (lower) $\beta_c$ means less tolerance for backorders.}
    \label{fig:problem-setup}
\end{figure}

Every period, the supplier makes two kinds of decisions, and this paper is concerned with both: (i) how much of each item to \emph{order} to replenish stock, and (ii) how to \emph{allocate} available on-hand stock across customers whenever it is insufficient to cover everyone immediately. Ordering follows a static base-stock policy with a fixed order-up-to level $S_i$ per item, set at the tactical level and held fixed in this paper. Allocation, by contrast, is a dynamic decision made every period: whenever on-hand stock cannot cover all outstanding demand for an item, the supplier must decide which customers to serve first. This is where the model's flexibility lies, and it is needed precisely because customers' SLAs are heterogeneous — some customers can tolerate more backorders than others, and an allocation policy that ignores this wastes stock protecting customers who do not need it, at the expense of customers who do.

Each customer $c$'s service level is expressed as a numeric backorder allowance $\beta_c \geq 0$: the maximum cumulative backorder units customer $c$ can accumulate across items within a review horizon of $T$ periods before a penalty applies. Concretely, if by the end of a review horizon customer $c$'s cumulative backorders exceed $\beta_c$ by some amount $x$, the supplier pays a penalty of $p_c \cdot x$, where $p_c > 0$ is the per-unit penalty cost (e.g., $\beta_1 = 6$ and $\beta_2 = 8$ units for two customers with different SLA strictness). Because review horizons repeat indefinitely over the supplier's operation, this is a recurring, long-run cost structure rather than a one-off constraint — the same recurring finite-horizon SLA structure studied for a single, aggregate service metric by \citet{temizoz2026ace}; here we instead track backorders per customer under heterogeneous, customer-specific allowances.

\section{Research Questions}

\begin{enumerate}
    \item[\textbf{RQ1}] \textit{Optimal Base-Stock Levels:} What base-stock levels $\mathbf{S}^* = \{S_i^*\}_{i \in I}$ minimize the supplier's long-run expected cost, subject to the customers' SLA backorder allowances?

    \item[\textbf{RQ2}] \textit{Optimal Allocation Policy:} Given fixed base-stock levels, what allocation policy $\pi: \mathcal{S} \to \mathcal{A}$ minimizes cost while ensuring that each customer's cumulative backorders respect their SLA allowance?
\end{enumerate}

\section{Model Formulation}

\subsection{Sets, Indices, and Time Notation}

\begin{itemize}
    \item $I$: set of items, indexed $i \in I$, $|I| \geq 1$
    \item $C$: set of customers, indexed $c \in C$, $|C| \geq 1$
    \item $t$: current period (absolute time index; $t = 1, 2, 3, \ldots$)
    \item $T$: review horizon length (number of periods per review horizon; e.g., $T = 20$)
    \item $r$: review horizon index — period $t$ belongs to horizon $r = r(t) = \lceil t/T \rceil$, which spans periods $(r-1)T+1$ through $rT$
    \item $L \geq 0$: replenishment lead time (periods between placing and receiving an order)
\end{itemize}

\textbf{Example:} With $T = 20$: periods 1--20 form horizon $r=1$; periods 21--40 form horizon $r=2$; period $t=35$ has $r(35) = \lceil 35/20 \rceil = 2$.

\subsection{Demand}

Demand for item $i$ from customer $c$ follows a known stochastic demand distribution (e.g., Poisson), independent across periods. Let $D_{c,i}(t) \in \mathbb{Z}_{\geq 0}$ denote the \emph{realized} demand of customer $c$ for item $i$ in period $t$ — a draw from that distribution, observed after the fact. Total demand is obtained by summing over customers (per item) or over items (per customer):
\begin{equation}
    D_i(t) = \sum_{c \in C} D_{c,i}(t), \qquad D_c(t) = \sum_{i \in I} D_{c,i}(t).
\end{equation}

\subsection{State Variables}

At period $t$, the state $s_t$ includes:

\subsubsection{On-Hand Inventory ($|I|$ features)}
\begin{equation}
    \text{OH}_i(t) \in \mathbb{Z}_{\geq 0}, \quad i \in I
\end{equation}
On-hand inventory of item $i$ at the \emph{start} of period $t$: before receiving the order $Q_i(t-L)$ placed $L$ periods ago, and before this period's order $Q_i(t)$ is placed. Once $Q_i(t-L)$ is received in Step 1 of the period workflow below, this becomes the before-allocation quantity $\text{OH}'_i(t)$, defined next.

\subsubsection{Pipeline Orders ($|I| \cdot L$ features)}
Orders placed in the past $L$ periods, currently in transit:
\begin{equation}
    Q_i(t-1), Q_i(t-2), \ldots, Q_i(t-L) \quad \text{for each item } i \in I
\end{equation}
These arrive in periods $t+1$ through $t+L$, respectively.

\subsubsection{Backorder per Customer-Item ($|C| \cdot |I|$ features)}
\begin{equation}
    \text{BO}_{c,i}(t) \in \mathbb{Z}_{\geq 0}, \quad c \in C, \, i \in I
\end{equation}
Units of item $i$ still owed to customer $c$ from past periods' unmet demand. A value of $\text{BO}_{c,i}(t) = 0$ simply means there are no backorders for that customer-item pair.

\subsubsection{Observed Demand, Current Period ($|C| \cdot |I|$ features)}
\begin{equation}
    D_{c,i}(t) \in \mathbb{Z}_{\geq 0}, \quad c \in C, \, i \in I
\end{equation}
as defined above.

\subsubsection{Cumulative Backorder Sum ($|C|$ features)}
Running sum of each customer's total backorder over the earlier periods of the current review horizon, \emph{not including period $t$ itself}:
\begin{equation}
    \overline{\text{BO}}_c(t) = \sum_{j=(r-1)T+1}^{t-1} \text{BO}_c(j), \qquad r = r(t),
\end{equation}
where $\text{BO}_c(j) = \sum_{i \in I} \text{BO}_{c,i}(j)$ is customer $c$'s total backorder at the \emph{beginning} of period $j$, i.e., what is owed from earlier periods' unmet demand. The sum is empty, and $\overline{\text{BO}}_c(t) = 0$, in the first period of every review horizon. It resets to 0 at the start of each new review horizon (see the update rule in Step 5 below).

\subsubsection{Time Remaining ($1$ feature)}
\begin{equation}
    T_{\text{rem}}(t) = r(t) \cdot T - t
\end{equation}

\subsubsection{SLA Backorder Allowance per Customer ($|C|$ features)}
\begin{equation}
    \beta_c \geq 0, \quad c \in C
\end{equation}
Maximum cumulative backorder units tolerated for customer $c$ per review horizon; a penalty is charged if $\overline{\text{BO}}_c(rT) + \text{BO}_c(rT) > \beta_c$ at horizon end, i.e., if the backorders at the beginning of the horizon's periods sum to more than $\beta_c$.

\textbf{Total state dimension:}
\begin{equation}
    |\mathcal{S}| = |I| + |I| \cdot L + |C|\cdot|I| + |C|\cdot|I| + |C| + 1 + |C|
\end{equation}

\subsection{On-Hand Inventory Before Allocation}

Once this period's arriving order has been received, on-hand inventory before allocation is
\begin{equation}
    \text{OH}'_i(t) = \text{OH}_i(t) + Q_i(t-L), \quad i \in I.
\end{equation}
This is the available supply for item $i$ this period — after the order $Q_i(t-L)$ has arrived but before this period's allocation is made. It is not an independent piece of state: it is fully determined by $\text{OH}_i(t)$ and the pipeline order $Q_i(t-L)$, both already part of the state. It is given its own symbol because it recurs throughout the rest of the model — in the inventory position below, in the allocation feasibility bound, and in Step 1 of the period workflow — and writing $\text{OH}_i(t) + Q_i(t-L)$ out in full each time would obscure that it is the same quantity every time.

\subsection{Inventory Position}

Ordering follows a base-stock policy, and a base-stock policy needs to know not just on-hand stock, but everything already committed to arrive and everything already owed to customers. This combined quantity is the inventory position — it would not be needed at all under a different ordering policy, but it is needed here specifically because ordering is base-stock:
\begin{equation}
    \text{IP}_i(t) = \text{OH}'_i(t) + \sum_{j=t-L+1}^{t-1} Q_i(j) - \sum_{c \in C} \text{BO}_{c,i}(t)
\end{equation}
i.e., on-hand inventory including this period's arriving order ($\text{OH}'_i(t)$), plus the remaining $L-1$ orders still in the pipeline (from the one placed two periods ago through the one placed yesterday, $Q_i(t-L+1), \ldots, Q_i(t-1)$), minus outstanding backorders across all customers. This is the same quantity as $\text{OH}_i(t) + \sum_{j=t-L}^{t-1} Q_i(j) - \sum_{c \in C} \text{BO}_{c,i}(t)$ — written via $\text{OH}'_i(t)$ because that quantity already has to be computed anyway (see above), not because a different set of orders is being counted.

\subsection{Decisions Each Period}

\subsubsection{Ordering Decision (static)}
\begin{equation}
    Q_i(t) = \max(0, S_i - \text{IP}_i(t)), \quad i \in I,
\end{equation}
where $S_i$ is the fixed base-stock level for item $i$. This order is placed in period $t$ and arrives in period $t+L$.

\subsubsection{Allocation Decision (dynamic)}
After demand is realized, an allocation policy $\pi(s_t)$ chooses, for each customer $c$ and item $i$:
\begin{equation}
    A_{c,i}(t) = \pi_{c,i}(s_t), \qquad \text{subject to} \qquad \sum_{c \in C} A_{c,i}(t) \leq \text{OH}'_i(t),
\end{equation}
i.e., total allocation of item $i$ cannot exceed the available supply $\text{OH}'_i(t)$ once this period's order has arrived. The policy can use any information in $s_t$, and may allocate less than what is available (a strategic hold) even when demand could be fully covered.

\newpage
\section{Period Workflow}

Each period $t$ executes the following steps in order.

\subsection{Step 1: Receive Orders}
The order $Q_i(t-L)$ placed $L$ periods ago arrives, becoming available alongside existing on-hand stock: $\text{OH}'_i(t) = \text{OH}_i(t) + Q_i(t-L)$, as defined above.

\subsection{Step 2: Compute Inventory Position and Place New Order}
$\text{IP}_i(t)$ is computed as defined above, and the new order $Q_i(t) = \max(0, S_i - \text{IP}_i(t))$ is placed (it will arrive in period $t+L$).

\subsection{Step 3: Observe Demand}
Demands $D_{c,i}(t)$ are realized for each customer-item pair; total item demand $D_i(t)$ is as defined above.

\subsection{Step 4: Allocate Using Policy}
For each item $i$, the allocation policy $\pi(s_t)$ chooses $A_{c,i}(t)$ for each customer $c$, subject to
\begin{equation}
    \sum_{c \in C} A_{c,i}(t) \leq \text{OH}'_i(t).
\end{equation}
The policy weighs backlog and new demand together — what is owed to customer $c$ for item $i$ is $\text{BO}_{c,i}(t) + D_{c,i}(t)$ — and decides how to split available stock across customers. If available supply covers everyone's backlog and new demand combined, the policy typically allocates in full, but it is not required to: the only hard constraint is the availability bound above.

\subsection{Step 5: Update State Variables}

\textbf{On-hand inventory:}
\begin{equation}
    \text{OH}_i(t+1) = \text{OH}'_i(t) - \sum_{c \in C} A_{c,i}(t)
\end{equation}
This is on-hand inventory at the start of period $t+1$ — before receiving $Q_i(t+1-L)$ and before placing $Q_i(t+1)$ — consistent with the definition of $\text{OH}_i(t)$ above. Since allocation is always bounded by available supply, on-hand inventory never goes negative; unmet demand is tracked separately as backorder.

\textbf{Backorder per customer-item:}
\begin{equation}
    \text{BO}_{c,i}(t+1) = \text{BO}_{c,i}(t) + D_{c,i}(t) - A_{c,i}(t)
\end{equation}
Unmet demand becomes backorder carried into the next period.

\textbf{Aggregate backorder per customer:}
\begin{equation}
    \text{BO}_c(t+1) = \sum_{i \in I} \text{BO}_{c,i}(t+1)
\end{equation}

\textbf{Cumulative backorder sum:}
\begin{equation}
    \overline{\text{BO}}_c(t+1) = \begin{cases}
        \overline{\text{BO}}_c(t) + \text{BO}_c(t) & \text{if } t < r(t) \cdot T \quad \text{(still inside the same horizon)} \\
        0 & \text{if } t = r(t) \cdot T \quad \text{(period $t+1$ starts a new horizon)}
    \end{cases}
\end{equation}
Individual backorders $\text{BO}_{c,i}(t+1)$ persist across horizon boundaries — only the cumulative sum $\overline{\text{BO}}_c$ restarts, at 0. Period $t+1$'s own backorder is not part of $\overline{\text{BO}}_c(t+1)$; it enters the sum at the next period, as $\text{BO}_c(t+1)$. In particular, the backorder created by the last period of a horizon is carried over and counts towards the next horizon, as the backorder at the beginning of that horizon's first period.

\textbf{Time remaining:}
\begin{equation}
    T_{\text{rem}}(t+1) = \begin{cases}
        T_{\text{rem}}(t) - 1 & \text{if } t < r(t) \cdot T \\
        T & \text{if } t = r(t) \cdot T
    \end{cases}
\end{equation}

\subsection{Step 6: Calculate Cost}

All costs are computed at period end, based on the state after allocation.

\textbf{Holding cost (every period):}
\begin{equation}
    C_{\text{hold}}(t) = \sum_{i \in I} h_i \cdot \text{OH}_i(t+1)
\end{equation}
charged on inventory remaining \emph{after} this period's allocation, i.e., what is left once this period's demand has been served.

\textbf{SLA penalty (at horizon end only):}
\begin{equation}
    C_{\text{penalty}}(t) = \begin{cases}
        \sum_{c \in C} p_c \cdot \max\!\Big(0,\; \underbrace{\overline{\text{BO}}_c(t)}_{\text{earlier periods this horizon}} + \underbrace{\text{BO}_c(t)}_{\text{beginning of this period}} - \; \beta_c\Big) & \text{if } T_{\text{rem}}(t) = 0 \text{ (end of horizon)} \\
        0 & \text{otherwise}
    \end{cases}
\end{equation}
$\overline{\text{BO}}_c(t)$ sums the backorder at the beginning of the horizon's periods up to but not including $t$; adding $\text{BO}_c(t)$ explicitly folds in the backorder at the beginning of the horizon's last period, so that the penalty covers the backorder at the beginning of every period of the horizon. The backorder created by the last period's own allocation, $\text{BO}_c(t+1)$, is not charged in this horizon; it is the beginning-of-period backorder of the next horizon's first period.

\textbf{Total period cost:}
\begin{equation}
    C(t) = C_{\text{hold}}(t) + C_{\text{penalty}}(t)
\end{equation}

\textbf{Parameters:} $h_i > 0$ holding cost per unit of item $i$; $p_c > 0$ penalty cost per unit of excess backorder for customer $c$; $\beta_c \geq 0$ backorder allowance for customer $c$ per review horizon.

\newpage
\section{Objective Function}
\label{sec:objective}

Because review horizons repeat indefinitely over the supplier's operation, the allocation policy is not chosen to minimize the cost of any single horizon — it is chosen to minimize the long-run cost \emph{per period}, taken over all horizons as time goes to infinity. This is the standard average-cost criterion for a recurring MDP \citep{temizoz2025dcl}, applied here across successive review horizons in the spirit of \citet{temizoz2026ace}:
\begin{equation}
    \min_{\pi} \; \limsup_{T' \to \infty} \frac{1}{T'} \, \mathbb{E}\left[\sum_{t=1}^{T'} C(t) \,\middle|\, s_0\right]
\end{equation}
where the expectation is over stochastic demand realizations $D_{c,i}(t)$, $\pi: \mathcal{S} \to \mathbb{R}^{|C| \times |I|}$ is the allocation policy determining $A_{c,i}(t)$ each period, and $C(t)$ is the total period cost defined in Step 6 above (holding cost every period, plus an SLA penalty at the end of each review horizon).

\textbf{Note:} The base-stock levels $S_i$ are fixed; only the allocation policy $\pi$ is subject to optimization here.

\section{Policy Learning Methodology (Future Work)}

The allocation policy $\pi(s_t)$ can be learned via various approaches (e.g., Deep Controlled Learning \citep{temizoz2025dcl}, reinforcement learning at industrial scale \citep{vanderhaar2026asml}, stochastic optimization). Such methodology is independent of the model formulation and will be addressed in separate research. This document provides the formal model; policy optimization is orthogonal to the model's correctness.

\section{Numerical Results}
\label{sec:results}
% TODO: the numbers below were produced with the code's earlier end-of-period backorder
% accounting (mdp.py _finalize_period). The model above now uses beginning-of-period
% backorders, so mdp.py must be updated and these experiments rerun before submission.

This section reports a first numerical study of the allocation decision (RQ2). The base-stock levels $S_i$ are held fixed and identical across items throughout, so any cost difference between policies is due to the allocation policy alone. Optimizing $\mathbf{S}$ (RQ1) is left for a later stage.

\subsection{Instance}

We use $|C| = 2$ customers and $|I| = 3$ items, with lead time $L = 2$ and the same review horizon $T = 20$ for both customers. Demand is Poisson, independent across periods and across customer-item pairs. Table~\ref{tab:instance} lists the parameters. The base-stock level $S_i = 6$ is the same for every item. It is deliberately tight relative to demand, so that stock is frequently scarce and the allocation decision matters.

\begin{table}[H]
    \centering
    \caption{Instance parameters.}
    \label{tab:instance}
    \begin{tabular}{lccc}
        \toprule
        & Item 1 & Item 2 & Item 3 \\
        \midrule
        Holding cost $h_i$ & 1.0 & 0.8 & 1.2 \\
        Base-stock level $S_i$ & 6 & 6 & 6 \\
        Demand rate $\lambda_{1,i}$ (customer 1) & 1.0 & 0.8 & 0.5 \\
        Demand rate $\lambda_{2,i}$ (customer 2) & 0.7 & 1.2 & 0.6 \\
        \bottomrule
    \end{tabular}

    \vspace{0.8em}

    \begin{tabular}{lcc}
        \toprule
        & Customer 1 & Customer 2 \\
        \midrule
        Backorder allowance $\beta_c$ & 6 & 8 \\
        Penalty cost $p_c$ & 50 & 60 \\
        Review horizon $T$ & 20 & 20 \\
        \bottomrule
    \end{tabular}
\end{table}

\subsection{Allocation Policies}

The allocation $A_{c,i}(t)$ is built one unit at a time: whenever stock of item $i$ is scarce, the policy is asked, once per available unit, to which customer that unit goes, or whether to hold it. The resulting allocation is the tally of these choices and covers every feasible allocation satisfying $\sum_{c} A_{c,i}(t) \leq \text{OH}'_i(t)$ and $A_{c,i}(t) \leq \text{BO}_{c,i}(t) + D_{c,i}(t)$. When available supply covers everything owed for an item, it is allocated in full without consulting the policy, which is never worse under the cost structure of Step~6. The paper prescribes no particular allocation policy, so we compare four rule-based benchmarks with a policy learned by Deep Controlled Learning (DCL) \citep{temizoz2025dcl}.

\begin{itemize}
    \item \textbf{FCFS:} serves the lowest-indexed customer that is still owed a unit.
    \item \textbf{SLA-gap:} serves the customer whose cumulative backorder $\overline{\text{BO}}_c$ exceeds its allowance $\beta_c$ by the most (or falls short of it by the least).
    \item \textbf{Cost-greedy:} serves the customer that would, if the unit were withheld, be projected to exceed $\beta_c$ with the highest penalty $p_c$.
    \item \textbf{Greedy-dynamic:} uses cost-greedy once any customer exceeds its allowance, SLA-gap once any customer has used more than $80\%$ of its allowance, and FCFS otherwise.
    \item \textbf{DCL:} a neural network that classifies the best allocation decision from the state. Training starts from greedy-dynamic, labels sampled states with the best action found by simulated lookahead, and repeats for three generations, each generation rolling out the previous one.
\end{itemize}

\subsection{Experimental Setup}

DCL used $n = 2000$ labeled samples per generation, $m = 100$ rollouts per candidate action, and a rollout horizon of $h = 50$ periods, for three generations. The policy network encodes each block of the state (inventory and pipeline, backorders, demand, cumulative backorder, time remaining, allowances) separately before a shared trunk of two layers of 128 units. It was trained with cross-entropy loss for up to 50 epochs with early stopping. Training took 242 seconds on a laptop CPU. The validation accuracy of the three generations was 81\%, 80\% and 83\%. This is a reduced budget: the default settings of the accompanying code ($n = 8000$, $m = 200$, $h = 100$) were not run to completion.

All policies were evaluated on the same 100 simulated trajectories, each of 1000 periods after a warm-up of 100 periods, using common random numbers. We report the long-run average cost per period, which is the objective of Section~\ref{sec:objective}, and the paired difference to a benchmark policy on the same demand sequences. The reported errors are standard errors of the mean over the 100 trajectories.

\subsection{Results}

\begin{table}[H]
    \centering
    \caption{Average cost per period ($S_i = 6$ for all items). Differences are paired against cost-greedy, the best rule-based policy. Standard errors are in parentheses.}
    \label{tab:results}
    \begin{tabular}{lccc}
        \toprule
        Policy & Cost per period & Difference to cost-greedy & Relative \\
        \midrule
        FCFS            & 65.17 (0.58) & $+12.16$ (0.09) & $+22.9\%$ \\
        SLA-gap         & 55.26 (0.57) & $+2.25$ (0.04)  & $+4.2\%$ \\
        Greedy-dynamic  & 53.46 (0.54) & $+0.45$ (0.02)  & $+0.8\%$ \\
        Cost-greedy     & 53.01 (0.54) & $0$             & $0\%$ \\
        \midrule
        DCL, generation 1 & 52.32 (0.53) & $-0.69$ (0.04) & $-1.3\%$ \\
        DCL, generation 2 & 52.45 (0.52) & $-0.56$ (0.04) & $-1.0\%$ \\
        DCL, generation 3 & 52.29 (0.53) & $-0.71$ (0.03) & $-1.3\%$ \\
        \bottomrule
    \end{tabular}
\end{table}

The ordering of the rule-based policies is clear: FCFS, which ignores the SLAs, costs about 23\% more than cost-greedy, while the SLA-aware rules are much closer to each other. All three DCL generations have lower average cost than every rule-based policy. The improvement over cost-greedy is small in relative terms, about 1\%, but it is many standard errors large because the comparison is paired. We do not observe a clear improvement from one DCL generation to the next: generation 2 is slightly worse than generations 1 and 3, and we have not tested whether the differences between generations are significant.

\paragraph{Limitations.} These results are for a single instance, a single base-stock level, one evaluation seed and a reduced training budget, and no repeated training runs were carried out. They show that a learned allocation can beat the best rule-based policy here, but they do not yet show how large the gain is in general, and the gain may depend on the base-stock level.

\section{Summary}

This model formalizes a multi-period, multi-customer, multi-item inventory system where:

\begin{itemize}
    \item \textbf{Orders} follow a static base-stock policy: $Q_i(t) = \max(0, S_i - \text{IP}_i(t))$

    \item \textbf{Allocations} are chosen dynamically by policy $\pi(s_t)$ to minimize long-run average cost

    \item \textbf{Backorders} are tracked per customer-item and accumulated within each review horizon, inclusive of the horizon's last period, to assess SLA compliance

    \item \textbf{Costs} consist of holding cost (every period, on post-allocation inventory) and SLA penalties (at horizon end, if cumulative backorders exceed the customer's allowance)

    \item \textbf{Research objective:} Find the optimal allocation policy $\pi^*$ that minimizes long-run average cost while respecting SLA constraints
\end{itemize}

%% ============================================================
%% Bibliography
%% ============================================================
\newpage
\bibliographystyle{plainnat}

\begin{thebibliography}{99}

\bibitem[Abbasi et~al.(2022)]{abbasi2022}
Abbasi, B., Fadaki, M., Hosseinifard, Z., Jahani, H., \& Thomas, D.~J. (2022).
Allocation policies to fulfil heterogeneous service requirements under resource pooling.
\textit{Decision Sciences}, 53(2), 277--319.

\bibitem[Boute et~al.(2022)]{boute2022}
Boute, R.~N., Gijsbrechts, J., van Jaarsveld, W., \& Vanvuchelen, N. (2022).
Deep reinforcement learning for inventory control: A roadmap.
\textit{European Journal of Operational Research}, 298(2), 401--412.

\bibitem[Sherbrooke(1968)]{sherbrooke1968}
Sherbrooke, C.~C. (1968).
METRIC: A multi-echelon technique for recoverable item control.
\textit{Operations Research}, 16(1), 122--141.

\bibitem[Sherbrooke(2004)]{sherbrooke2004}
Sherbrooke, C.~C. (2004).
\textit{Optimal Inventory Modeling of Systems: Multi-Echelon Techniques}.
Kluwer Academic Publishers.

\bibitem[Tan et~al.(2017)]{tan2017}
Tan, Y.~R., Paul, A.~A., Deng, Q., \& Wei, L. (2017).
Mitigating inventory overstocking: Optimal order-up-to level to achieve a target fill rate over a finite horizon.
\textit{Production and Operations Management}, 26(11), 1971--1988.

\bibitem[Temizoz et~al.(2025)]{temizoz2025dcl}
Temizoz, T., Imdahl, C., Dijkman, R., Lamghari-Idrissi, D., \& van Jaarsveld, W. (2025).
Deep Controlled Learning for Inventory Control.
\textit{European Journal of Operational Research}, 324(1), 104--117.

\bibitem[Temizoz et~al.(2026)]{temizoz2026ace}
Temizoz, T., Imdahl, C., Dijkman, R., Lamghari-Idrissi, D., \& van Jaarsveld, W. (2026).
How to ace your next service level contract review?
SSRN Working Paper No.~6075348. \url{https://papers.ssrn.com/sol3/papers.cfm?abstract_id=6075348}

\bibitem[van der Haar et~al.(2026)]{vanderhaar2026asml}
van der Haar, J.~F., van Jaarsveld, W., Basten, R.~J.~I., \& Boute, R.~N. (2026).
Industrializing deep reinforcement learning for operational spare parts inventory management.
\textit{European Journal of Operational Research}, 334, 128--140.

\end{thebibliography}

\end{document}
